\documentclass[12pt]{article}
\usepackage[a4paper,margin=1in]{geometry}
\usepackage[T1]{fontenc}
\usepackage[utf8]{inputenc}
\usepackage{hyperref}
\usepackage{longtable}
\usepackage{tabularx}
\usepackage{array}
\usepackage{enumitem}

\hypersetup{
  colorlinks=true,
  linkcolor=black,
  urlcolor=blue
}

\setlist[itemize]{leftmargin=1.6em}
\setlist[enumerate]{leftmargin=1.8em}
\renewcommand{\arraystretch}{1.2}

\title{QuickCart Blackbox Testing Report}
\author{}
\date{}

\begin{document}
\maketitle

\section{Introduction}
In this part of the assignment, I tested the QuickCart server as a closed system. I did not use any internal backend source code while designing the tests. Instead, I treated the API documentation and the live HTTP responses as the only specification, and I used the published admin inspection endpoints whenever I needed to verify state changes from the outside.

My goal was to build a practical blackbox test suite that covers normal user flows, invalid inputs, boundary values, and state transitions across the main API areas: profile, addresses, products, cart, coupons, checkout, wallet, loyalty, orders, reviews, and support tickets.

\section{Blackbox Testing Approach}
I followed three main blackbox techniques while designing the suite:

\begin{itemize}
  \item \textbf{Equivalence partitioning:} I grouped requests into valid and invalid classes, for example valid vs invalid headers, valid vs invalid profile updates, valid vs invalid address labels, and allowed vs disallowed support ticket transitions.
  \item \textbf{Boundary value analysis:} I tested the edge of documented limits such as 0 quantity, phone numbers that were too short or too long, 6-digit pincodes, loyalty redemption below 1, wallet add limits, comment length 0 and 201, and COD totals above 5000.
  \item \textbf{State transition testing:} I checked flows where one action should change later behavior, such as default address switching, cart accumulation, coupon application and removal, checkout creating orders, order cancellation affecting stock, review submission affecting average rating, and support tickets moving from OPEN to IN\_PROGRESS to CLOSED.
\end{itemize}

I also used the admin inspection endpoints as part of the external interface. For example, I used them to confirm whether inactive products were hidden from the public list, whether product prices matched the stored values, whether support ticket messages were saved exactly, and whether stock was restored after cancellation.

\section{Environment Setup}
I ran the system from the provided Docker image. On this machine, the working image was the ARM64 image file \texttt{quickcart\_image.tar}. The minimal reproducible setup is:

\begin{verbatim}
docker load -i quickcart_image.tar
docker run --rm --name quickcart-blackbox -p 8080:8080 quickcart:latest
python3 -m pytest blackbox/tests -v
\end{verbatim}

I used the following defaults in the test suite:

\begin{itemize}
  \item Base URL: \texttt{http://127.0.0.1:8080}
  \item Roll number header: \texttt{23110001}
\end{itemize}

No tests were blocked by the environment once the container was running.

\section{Test Scenarios Covered}
I organized the suite feature by feature so the test names stay readable and each file focuses on one area of observable behavior.

\subsection{Header Validation}
I tested missing \texttt{X-Roll-Number}, malformed \texttt{X-Roll-Number}, missing \texttt{X-User-ID}, non-positive \texttt{X-User-ID}, non-existent user IDs, and the rule that admin endpoints do not require \texttt{X-User-ID}. These cases were important because every other endpoint depends on correct header validation.

\subsection{Profile}
I tested profile retrieval, valid profile updates, invalid name boundaries, invalid phone length, and non-digit phone input. These checks matter because profile updates are simple to use but easy to validate incorrectly, especially around boundary cases.

\subsection{Addresses}
I tested valid address creation, invalid label/street/city/pincode inputs, update restrictions, deleting a missing address, and the one-default-address rule. I focused on these because address handling is stateful and the default flag can easily become inconsistent across records.

\subsection{Products}
I tested public product listing, product lookup by ID, invalid product IDs, category filtering, name search, price sorting in both directions, and price consistency against the admin state. These tests mattered because the product catalog is the starting point for later cart and checkout flows.

\subsection{Cart and Coupons}
I tested empty and non-empty cart views, adding valid items, accumulating quantity when the same product is added twice, updating quantity, removing items, invalid quantities, missing products, subtotal math, total math, valid fixed and percentage coupons, minimum-cart checks, coupon removal, and expired coupon handling. I gave extra attention here because the cart is the most arithmetic-heavy part of the API and small mistakes here affect checkout and invoices later.

\subsection{Checkout, Wallet, Loyalty, and Orders}
I tested invalid payment methods, empty-cart checkout, COD over 5000, CARD/COD/WALLET payment statuses, wallet balance lookup, wallet add and pay behavior, invalid wallet amounts, loyalty balance lookup, valid and invalid loyalty redemption, order listing, order detail lookup, cancellation rules, invoice math, and stock restoration after cancellation. These scenarios matter because they connect multiple API areas and reveal whether the system preserves correct financial state across actions.

\subsection{Reviews and Support Tickets}
I tested the no-review average boundary, valid review submission, average rating as a decimal, invalid rating values, invalid comment lengths, valid support ticket creation, exact message storage, subject/message length validation, and allowed vs invalid ticket status transitions. These checks were useful because both features combine validation with visible persistent state.

\section{Test Execution Summary}
The final automated run was:

\begin{verbatim}
python3 -m pytest blackbox/tests -v
\end{verbatim}

Final result:

\begin{itemize}
  \item Total tests collected: 77
  \item Passed: 66
  \item Expected failures (\texttt{xfail}): 11
  \item Unexpected failures: 0
\end{itemize}

I kept the known server defects as \texttt{xfail} cases because the repository only provides a Dockerized server artifact for blackbox testing, not editable backend source code. This let me keep the suite runnable while still preserving the mismatches between the documented behavior and the observed behavior.

\section{Failures and Defects Observed}
During live execution, I found the following externally visible defects:

\begin{longtable}{|>{\raggedright\arraybackslash}p{0.18\textwidth}|>{\raggedright\arraybackslash}p{0.24\textwidth}|>{\raggedright\arraybackslash}p{0.24\textwidth}|>{\raggedright\arraybackslash}p{0.24\textwidth}|}
\hline
\textbf{Endpoint / Area} & \textbf{Expected Behavior} & \textbf{Actual Behavior} & \textbf{Impact} \\
\hline
\endfirsthead
\hline
\textbf{Endpoint / Area} & \textbf{Expected Behavior} & \textbf{Actual Behavior} & \textbf{Impact} \\
\hline
\endhead
\texttt{PUT /api/v1/profile} & Phone numbers should be exactly 10 digits and digits only. & Alphabetic phone content such as \texttt{12345abcde} is accepted. & Invalid user data can be stored in the profile. \\
\hline
\texttt{POST /api/v1/addresses} & Adding a new default address should unset the old default first. & Multiple addresses remain marked as default at the same time. & Default address selection becomes ambiguous. \\
\hline
\texttt{POST /api/v1/cart/add} & Quantity 0 should be rejected with a 400 error. & Quantity 0 is accepted. & Invalid cart state can enter the system. \\
\hline
\texttt{GET /api/v1/cart} subtotal & Item subtotal should equal quantity times unit price. & Some public cart flows return impossible subtotals, including negative values. & Cart math becomes unreliable. \\
\hline
\texttt{GET /api/v1/cart} total & Cart total should equal the sum of all item subtotals. & The total can stay at 0 or otherwise fail to match the item lines. & Later financial calculations become hard to trust. \\
\hline
\texttt{POST /api/v1/coupon/apply} & Expired coupons should be rejected. & Expired coupons are still accepted when the cart meets the minimum amount. & Users can receive discounts that should not be available. \\
\hline
\texttt{POST /api/v1/checkout} with empty cart & Empty-cart checkout should fail. & The API creates a zero-total order instead of rejecting the request. & Invalid orders can be created from no items. \\
\hline
\texttt{POST /api/v1/wallet/pay} & The exact requested amount should be deducted. & The wallet balance drops by more than the requested amount. & Wallet balances become inaccurate. \\
\hline
\texttt{POST /api/v1/orders/\{id\}/cancel} & Cancelling an order should restore the ordered stock. & The order is marked cancelled, but stock does not return. & Inventory stays lower than it should. \\
\hline
\texttt{GET /api/v1/orders/\{id\}/invoice} & Invoice total should match subtotal plus one GST charge. & The total behaves like GST is added twice. & Invoice totals do not match the actual order amount. \\
\hline
\texttt{POST /api/v1/products/\{id\}/reviews} & Ratings outside 1 to 5 should be rejected. & Rating 6 is accepted. & Review data and average ratings can become invalid. \\
\hline
\end{longtable}

I did not apply backend fixes for these issues because the blackbox setup only exposes the running server image, not an editable server codebase inside the repository. Instead, I documented each defect and preserved it in the automated suite as an expected-failure case.

\section{How to Run the Code and Tests}
The blackbox deliverables can be run with:

\begin{verbatim}
docker load -i quickcart_image.tar
docker run --rm --name quickcart-blackbox -p 8080:8080 quickcart:latest
python3 -m pytest blackbox/tests -v
\end{verbatim}

Optional environment variables:

\begin{itemize}
  \item \texttt{QUICKCART\_BASE\_URL}
  \item \texttt{ROLL\_NUMBER}
\end{itemize}

\section{Final Reflection}
This part was a good example of why blackbox testing needs both documentation and live probing. The API documentation gave me the intended rules, but the live server responses showed where those rules were not actually enforced. I found that the most useful tests were the ones that combined boundary values with multi-step state changes, especially in cart, checkout, wallet, and order flows.

I also found that the admin inspection endpoints were very helpful because they let me verify state transitions without breaking the blackbox approach. They are part of the public interface, so they gave me a clean way to check prices, stock changes, coupon data, and stored support ticket messages from the outside.

Overall, the final suite is runnable, covers all the required API areas, and clearly separates normal passing behavior from real defects in the provided QuickCart image.

\end{document}
